\noindent The pinned model versions: every scheme's primary scorer is
\texttt{gpt-4o-mini}, pinned to the dated snapshot
\texttt{gpt-4o-mini-2024-07-18}; the whole-grid sensitivity re-score uses
\texttt{gpt-4o}, pinned to \texttt{gpt-4o-2024-08-06}
(\texttt{model\_pins.py}, verified against the provider's model list before
scoring).\medskip

\subsection*{1. Inspection reports --- Warmth}
What is measured: how the report describes the relationships between staff and pupils. Score the relational quality described, never the pastoral provision listed.

\begin{enumerate}
  \item Relationships described as poor: pupils unknown to staff, uncared for, or unsafe because of how staff relate to them.
  \item Any credible negative relational signal --- pupils not always heard, bullying unresolved, staff not accessible --- even alongside positives.
  \item DEFAULT. A safe and generally positive school, described through labels, provision or outcomes ("caring school", "pupils feel safe", "pastoral team is effective") without describing how staff and pupils actually relate.
  \item The report describes the relationships themselves: staff know pupils well, warm interactions observed, pupils feel valued by staff as individuals.
  \item Relational warmth described as a defining feature of the school, in emphatic terms, across several passages.
\end{enumerate}
Rules: if the report does not discuss relationships, score 3 --- silence is not coldness. The 3/4 line: does the sentence describe how staff and pupils relate, or the school's character and what provision exists? In doubt between two bands, take the lower.

\subsection*{2. Inspection reports --- Strictness}
What is measured: what the report says about conduct standards and how consistently the adults enforce them.

\begin{enumerate}
  \item Disruption is a persistent, documented problem; rules unclear or not applied.
  \item Standards exist but enforcement is patchy; behaviour is named as an area for improvement, or inspectors saw low-level disruption.
  \item DEFAULT. Clear expectations, generally maintained; behaviour mentioned but unremarkable.
  \item The report as a whole presents conduct as a strength, in lessons and around the building.
  \item Conduct is exceptional and unmistakably a defining feature of the school.
\end{enumerate}
Rules: if behaviour is simply not discussed, score 3 --- a 1 or 2 needs positive evidence of a problem. Pupils' opinions about the rules are warmth evidence, never strictness evidence: strictness is what the adults do. In doubt between two bands, take the lower.

\subsection*{3. Inspection reports --- Teaching}
What is measured: how the report describes classroom teaching --- explanations, subject knowledge, and the checking of understanding. Curriculum-design language (intent, sequencing) is background, not evidence.

\begin{enumerate}
  \item Teaching named as a serious failure and a cause of poor outcomes.
  \item Teaching inconsistent ("some teachers... not always") and named as requiring improvement.
  \item DEFAULT. Competent and unremarkable, or classroom delivery not discussed at all.
  \item Teaching described positively and with specifics across most of the school ("typically", "most teachers").
  \item Teaching described as expert and a defining strength ("consistently", "across all subjects").
\end{enumerate}
Rules: score how teachers teach, not what the curriculum intends. The consistency words carry the ladder: "consistently" points to 5, "typically" to 4, "some" to 2. In doubt, or with no delivery evidence, score 3.

\subsection*{4. Behaviour policies --- Strictness}
What is measured: how far the policy removes staff discretion and prescribes pupil conduct. Score the machinery the policy describes, not how severe its language sounds.

\begin{enumerate}
  \item No consequence language at all --- the policy is aspirational or entirely about rewards.
  \item Consequence language is present, but no day-to-day system a new member of staff could act on tomorrow.
  \item DEFAULT. A working day-to-day system, applied by staff judgement from a menu of options. The standard English secondary policy.
  \item A named tracking or tier system: a pupil's position in the system, not staff judgement, decides which consequence follows.
  \item As 4, and the policy prescribes a default deportment norm --- a SLANT-type routine or a similar required way of sitting, attending or moving.
\end{enumerate}
Rules: the test for 3 --- could a new teacher read this policy and know what to do tomorrow when a lesson is disrupted? Harsh penalties at staff discretion are LESS strict on this construct than mild ones applied automatically.

\subsection*{5. Behaviour policies --- Warmth}
What is measured: what the policy tells staff about how to relate to pupils. Rewards and pastoral structures are provision, not warmth.

\begin{enumerate}
  \item No positive or relational content beyond the sanctions system.
  \item Positive values asserted ("every child matters") with no described mechanism and no guidance to staff.
  \item DEFAULT. A positive mechanism described --- named rewards with the occasions they are given on, or named pastoral routes --- but no guidance to staff on how to relate to pupils. Most policies.
  \item Explicit operational guidance to staff on relating to pupils as individuals: greet at the door, learn names, how to hold a repair conversation.
  \item As 4, and relational quality is the organising principle: consequences framed as relationship-preserving, and the policy says the relationship survives the sanction.
\end{enumerate}
Rules: a scheduled restorative meeting is machinery; only telling staff how to conduct it is guidance. If in doubt whether a sentence is guidance, it is not --- a rich reward economy cannot pass 3.

\subsection*{6. Websites --- Strictness}
What is measured: whether the site's high-expectations language is about behaviour at all, and how far the school goes in claiming strictness of itself.

\begin{enumerate}
  \item No mention of high standards or high expectations anywhere --- even a cursory mention moves it up.
  \item High standards or expectations claimed, but generic --- not tied to behaviour or conduct.
  \item DEFAULT. High expectations explicitly about behaviour or conduct --- or the school asserting it is demanding, making no apologies.
  \item Behaviour treated as central: expected at all times, pupils required to become self-disciplined, rewards-and-sanctions machinery on the site, or behaviour named as the foundation of success.
  \item The school states that it is strict: strict routines and protocols, zero tolerance, no excuses, the highest possible standards.
\end{enumerate}
Rules: aspiration language alone ("achieve their full potential") does not reach band 2. It is the explicit statement of strictness that earns the 5, not the severity of the machinery.

\subsection*{7. Websites --- Warmth}
What is measured: how far the site presents care for pupils and staff--pupil relationships as what the school is about, in the school's own words. The axis: absent --- courtesy --- claimed --- stressed --- central.

\begin{enumerate}
  \item No care or support language at all. The school presents itself entirely through results, facilities, opportunities and ambition.
  \item The courtesy register only: generic supportive-atmosphere phrases ("a supportive learning environment", "care and support for each individual") appearing in passing, typically once inside the head's welcome, and not returned to.
  \item DEFAULT. Care claimed as a feature of the school: pastoral care named, support for individual pupils emphasised, or a caring atmosphere claimed prominently or in more than one place.
  \item Care or relationships stressed as important and returned to: pupils described as known and valued as individuals, care presented as part of what the school is for, across more than one page or passage.
  \item Staff--pupil relationships explicitly named and presented as central to what the school is --- and more than once. Rare by design.
\end{enumerate}
Rules: the 2/3 line is courtesy versus feature --- if the care language lives only in the head's welcome paragraph, score 2. Band 5 needs the relationships themselves named (not just "caring" or "supportive"), plus centrality, plus recurrence. Quoted Ofsted praise does not count; the school's own and its trust's words do. In doubt between two bands, take the lower.

